% Phase-two research fragment. Uses the theorem environments and packages
% of research-progress.tex. The first-moment kernel estimate is restated here.
\section{A coefficient-one random-host lower bound}
\label{sec:rhost-phase-two}

\paragraph{Status.}
This section improves the lower coefficient from $2/3$ to $1$.
It uses a weighted pseudoforest argument and the same random-host kernel count.
The proof is not formalised in Lean.
We have not established publication priority.

All logarithms are natural, except $\log_2$.
A pseudoforest is a graph whose components contain at most one cycle each.
All graphs in this section are finite and simple.
Isolated vertices may be retained in every spanning subgraph.

\begin{lemma}[Cycles and pseudoforests]
\label{lem:rhost-pseudoforest-cycles}
Let $W$ be a graph, let $P\subseteq W$ be a spanning pseudoforest,
and let $\mathcal C$ be a family of $k$ distinct cycles of $W$.
Then
\[
 \sum_{C\in\mathcal C}|E(C)\setminus E(P)|
 \ge 2\bigl(k-\beta(W)\bigr).
\]
\end{lemma}

\begin{proof}
Let $u$ be the number of unicyclic components of $P$.
Let $a$ count omitted edges whose endpoints lie in the same component of $P$.
Let $b$ count the remaining omitted edges.
If $c_P$ and $c_W$ are the component counts, then $b\ge c_P-c_W$.
Consequently,
\[
 \beta(W)=u+a+b-c_P+c_W\ge u+a.
\]

Write $n_j$ for the number of selected cycles using exactly $j$ omitted edges.
A unicyclic component has only one cycle, so $n_0\le u$.
A tree has at most one simple path between two fixed vertices.
A unicyclic graph has at most two.
Thus each of the $a$ internal omitted edges belongs to at most two cycles counted by $n_1$.
An omitted edge between different components belongs to none.
It follows that
\[
 2n_0+n_1\le2u+2a\le2\beta(W).
\]
Therefore
\[
 \sum_{C\in\mathcal C}|E(C)\setminus E(P)|
 =\sum_{j\ge0}jn_j
 \ge2k-2n_0-n_1
 \ge2(k-\beta(W)).
\]
\end{proof}

\begin{lemma}[Balanced pseudoforest cover]
\label{lem:rhost-balanced-cover}
Let $R\ge2$ be an integer.
Suppose every nonempty edge set $A\subseteq E(W)$ satisfies
\[
 (R-1)|A|\le R|V(A)|,
\]
where $V(A)$ is the set of endpoints of edges in $A$.
Then there are $R$ spanning pseudoforests $P_1,\ldots,P_R$ of $W$
such that each edge of $W$ is absent from exactly one $P_i$.
\end{lemma}

\begin{proof}
Form a bipartite graph.
Its left side has $R-1$ copies of every edge of $W$.
Its right side has $R$ copies of every vertex.
Each edge-copy is adjacent to all copies of both endpoints.
For a set $X$ of left vertices, let $A$ be the original edges represented in $X$.
Then
\[
 |X|\le(R-1)|A|\le R|V(A)|=|N(X)|.
\]
Hall's theorem gives a matching that covers every left vertex.

Forget the copy labels after recording the matched incidences.
This gives a bipartite multigraph with sides $E(W)$ and $V(W)$.
Its edge-vertices have degree $R-1$, and its original vertices have degree at most $R$.
Colour its edges properly with $R$ colours.
For completeness, such a colouring follows again from Hall's theorem.
Add dummy vertices and edges to make both sides equally large and every degree equal to $R$.
A regular bipartite multigraph has a perfect matching by Hall's theorem.
Remove one perfect matching at a time to obtain the $R$ colours.

For colour $i$, retain every original edge incident with an incidence of that colour.
Each retained edge has an assigned endpoint, and no two retained edges have the same assigned endpoint.
Every component therefore has at most as many edges as vertices.
The retained graph is a pseudoforest.
The $R-1$ incidences at each edge-vertex have distinct colours.
Thus each original edge is absent from exactly one colour.
\end{proof}

\begin{proposition}[Weighted cycle-length bound]
\label{prop:rhost-weighted-cycles}
Under the hypotheses of Lemma~\ref{lem:rhost-balanced-cover}, every family
$\mathcal C$ of distinct cycles satisfies
\[
 \boxed{\quad
 \beta(W)\ge\sum_{C\in\mathcal C}
 \left(1-\frac{|E(C)|}{2R}\right).
 \quad}
\]

In particular, one may sum over all cycles of length less than $2R$.
\end{proposition}

\begin{proof}
Apply Lemma~\ref{lem:rhost-pseudoforest-cycles} to every $P_i$ from
Lemma~\ref{lem:rhost-balanced-cover}.
Every edge is omitted exactly once, so
\[
 \sum_{i=1}^{R}\sum_{C\in\mathcal C}|E(C)\setminus E(P_i)|
 =\sum_{C\in\mathcal C}|E(C)|.
\]
The same sum is at least $2R(|\mathcal C|-\beta(W))$.
Rearrange this inequality.
\end{proof}

This deterministic inequality allows dependencies between the selected cycles.
It therefore avoids the theta obstruction to their complete independence.

\begin{lemma}[Local sparsity of the random host]
\label{lem:rhost-phase-two-sparsity}
Let $G\sim G(n,D/n)$, with $D\ge2$.
Put
\[
 R=\left\lfloor\frac{\ln n-40\ln\ln n}{\ln D}\right\rfloor,
 \qquad S=\lceil8(\ln n)^2\rceil.
\]
Assume $R\ge3$.
With probability $1-O((\ln n)^{-17})$, every subgraph $W$ with at most $S$ edges
satisfies the hypothesis of Lemma~\ref{lem:rhost-balanced-cover}.
\end{lemma}

\begin{proof}
The kernel enumeration in Lemma~\ref{lem:rhost-kernel} applies with this $S$ and $R$.
It bounds the probability of a connected subgraph $F$ satisfying
\[
 \delta(F)\ge2,\qquad q=|E(F)|-|V(F)|\ge1,
 \qquad |E(F)|\le\min\{S,Rq\}
\]
by
\[
 S\sum_{q=1}^{S}
 \left(\frac{(2S)^{10}D^R}{n}\right)^q
 =O((\ln n)^{-17}).
\]
Indeed, $D^R\le n/(\ln n)^{40}$ and $(2S)^{10}\le(\ln n)^{21}$ for large $n$.
Assume no such $F$ occurs.

Consider a connected component of an edge set $A\subseteq E(W)$.
Write $e$ and $v$ for its edge and vertex counts.
If $e\le v$, then $(R-1)e\le Rv$ immediately.
If $q=e-v\ge1$, repeatedly remove vertices of degree one.
The remaining nonempty core is connected, has minimum degree at least two, and retains the same $q$.
If the original component had $e\le Rq$, then its core would be forbidden.
Thus $e>Rq$, which implies $(R-1)e<Rv$.
Sum these inequalities over the components of $A$.
\end{proof}

\begin{theorem}[Coefficient-one lower bound]
\label{thm:rhost-coefficient-one}
Let $G\sim G(n,D/n)$, and let $R$ be as in Lemma~\ref{lem:rhost-phase-two-sparsity}.
Assume $D\ge2$ and $R\ge3$.
With probability $1-O((\ln n)^{-17})$, every spanning pancyclic subgraph $H$ satisfies
\[
 \boxed{\quad
 |E(H)|-n
 \ge\log_2(n-2)+R-3-\log_2(2R^2-R-3).
 \quad}
\]
In particular, if $D\sim\ln n$, then with high probability
\[
 \operatorname{Pex}(G)
 \ge\log_2 n+
 \frac{\ln n}{\ln D}
 -2\log_2\!\left(\frac{\ln n}{\ln D}\right)-O(1).
\]
Consequently,
\[
 \operatorname{Pex}(G)
 \ge\log_2 n+(1-o(1))\frac{\ln n}{\ln\ln n}.
\]
\end{theorem}

\begin{proof}
Choose in $H$ one cycle of every length $3,\ldots,2R-1$.
Let $W$ be their union.
Its number of edges is at most
\[
 M=\sum_{j=3}^{2R-1}j=2R^2-R-3\le S.
\]
On the event in Lemma~\ref{lem:rhost-phase-two-sparsity},
Proposition~\ref{prop:rhost-weighted-cycles} gives
\[
 \beta(W)\ge(2R-3)-\frac{2R^2-R-3}{2R}
 =R-\frac52+\frac{3}{2R}.
\]
Since $\beta(W)$ is an integer, $\beta(W)\ge R-2$.

Put $r=\beta(H)$.
The outside-part length count yields
\[
 n-2\le |E(W)|2^{r-\beta(W)}\le M2^{r-R+2}.
\]
The graph $H$ is connected, so $|E(H)|-n=r-1$.
Taking binary logarithms proves the finite bound.
If $D\sim\ln n$, then $R=\ln n/\ln D-O(1)$.
Also, $\log_2(2R^2-R-3)=2\log_2(\ln n/\ln D)+O(1)$.
\end{proof}

At the explicit probability $p_*$ from Theorem~\ref{thm:rhost-scale},
this improves the proved interval to
\[
 \log_2n+(1-o(1))\frac{\ln n}{\ln\ln n}
 \le\operatorname{Pex}(G(n,p_*))
 \le\log_2n+(2+o(1))\frac{\ln n}{\ln\ln n}.
\]
The upper bound and its probability range remain unchanged.

\paragraph{Scope of the improvement.}
The weighted inequality is deterministic and has no independence hypothesis on the cycles.
The random input only ensures the local edge-density condition.
The outside-part count then adds the rank cost of the short lengths to $\log_2n$.
This proof does not determine the sharp correction constant or extend the upper bound below $p_*$.

\subsection{Replacing individual short cycles by theta blocks}
\label{subsec:rhost-theta-upper}

\begin{lemma}[A theta block with prescribed lengths]
\label{lem:rhost-theta-block}
Let $s,x,y$ be integers at least three, with $x<y$.
Suppose
\[
 s+x+y\equiv0\pmod2,\qquad y-x<s<x+y.
\]
Set
\[
 \alpha=\frac{s+x-y}{2},\qquad
 \gamma=\frac{s+y-x}{2},\qquad
 z=\frac{x+y-s}{2}.
\]
An $s$-cycle with an additional internally disjoint path of length $z$
between two vertices at cyclic distance $\alpha$ contains cycles of lengths $s,x,y$.
Its rank is two.
It has a spanning path whose endpoints can be specified before the added path is embedded.
\end{lemma}

\begin{proof}
The three internally disjoint paths between the branch vertices have lengths
$\alpha,\gamma,z$.
Their pairwise sums are $s,x,y$.
Write the base cycle as $v_0v_1\cdots v_{s-1}v_0$, with branch vertices
$u=v_0$ and $v=v_\alpha$.
Delete the base edges $v_0v_1$ and $v_\alpha v_{\alpha+1}$.
Retain the full added path between $u$ and $v$.
The result is a spanning path with endpoints $v_1$ and $v_{\alpha+1}$.
This also applies when $\alpha=1$.
\end{proof}

\begin{lemma}[Packing the short length interval]
\label{lem:rhost-theta-interval}
Let $d,b$ be positive integers with $b=o(d)$ along the sequence under consideration.
Put
\[
 T=2(d+b+1),\qquad H=d+5,\qquad
 J=\left\lfloor\frac{T-H+1}{8}\right\rfloor.
\]
For all sufficiently large $d$, one can cover every length $3,\ldots,T$
with disjoint abstract blocks of total rank
\[
 T-2-4J\le\frac32d+b+5.
\]
Each block is either a cycle or a theta graph.
Every theta block has a base cycle of length at most $4J+2$,
and its additional path has length at least $d+4$ and at most $O(d)$.
\end{lemma}

\begin{proof}
For $0\le j<J$, write $h=H+8j$ and use the following four triples $(s,x,y)$:
\[
 \begin{array}{c|c|c|c}
 s&x&y&z=(x+y-s)/2\\ \hline
 4j+3&h+4&h+5&H+6j+3\\
 4j+4&h&h+2&H+6j-1\\
 4j+5&h+6&h+7&H+6j+4\\
 4j+6&h+1&h+3&H+6j-1
 \end{array}
\]
Each row satisfies Lemma~\ref{lem:rhost-theta-block}.
These blocks cover exactly the two disjoint sets
\[
 \{3,\ldots,4J+2\}\quad\text{and}\quad\{H,\ldots,H+8J-1\}.
\]
Their disjointness follows from $b=o(d)$.
Use one cycle for every remaining length in $\{3,\ldots,T\}$.
There are $4J$ theta blocks and $T-2-12J$ single-cycle blocks.
Their total rank is $8J+(T-2-12J)=T-2-4J$.
The claimed estimate follows from the definition of $J$.
\end{proof}

\begin{theorem}[An upper coefficient of three halves]
\label{thm:rhost-three-halves}
For the same $p_*$ as in Theorem~\ref{thm:rhost-scale}, and every $p\ge p_*$,
with high probability
\[
 \operatorname{Pex}(G(n,p))
 \le\log_2n+\left(\frac32+o(1)\right)\frac{\ln n}{\ln\ln n}.
\]
\end{theorem}

\begin{proof}
Keep $p_1,\ldots,p_5,\beta,a,d,K_0,K,L_0$ from the audited construction above.
Keep the two-digit gadget with $b=\lceil\sqrt{2\ln n}\rceil$.
Its shortest represented length is $T+1$, where $T=2(d+b+1)$.
We replace the individual short chords by the blocks in Lemma~\ref{lem:rhost-theta-interval}.
Write $A=4J$ for the number of theta blocks and $B=T-2-8J$ for the number of all short blocks.
Their total rank is $A+B=T-2-4J$.

\emph{Initial components.}
Split the $G_1$ exposure into four independent subexposures with equal edge probability
$q_1=1-(1-p_1)^{1/4}$.
Then $nq_1\to\infty$.
Use the first subexposure for the original binary base cycles $C_0,\ldots,C_{K_0}$.
Use the second for one base cycle of each short block.
Use the third for the two-digit base cycle $C_{\rm short}$.
The cycle-length input used in the original construction applies to each subexposure.

All requested lengths are at most the largest original binary base-cycle length, for sufficiently large $n$.
The original short-cycle separation estimate therefore makes all distinct selected cycles vertex-disjoint.
Cycles selected from different subexposures are also distinct with high probability.
Indeed, each pairwise intersection has edge probability $q_1^2$ and $nq_1^2=o(1)$.
Its expected total number of cycles is $o(1)$.
Finally, choose an edge $f$ in the fourth subexposure outside all selected vertices.
This is possible with high probability, conditionally on the earlier selections.

For each theta base cycle, fix its branch vertices and its spanning-path endpoints from
Lemma~\ref{lem:rhost-theta-block}.
For each single-cycle block, delete a designated edge to specify its spanning-path endpoints.
All these endpoints are now fixed, before $G_2$ is exposed.

\emph{The main cycle and the chain of short blocks.}
Build the binary base cycle $C_*$ through the designated edges of
$C_0,\ldots,C_{K_0}$ and through the extra edge $f$.
Use $K_0+2$ internally disjoint $G_2$ paths, each of length $d+2$.
Thus
\[
 |C_*|=(K_0+2)(d+3)=(1+o(1))\log_2n.
\]
Reserve $f$ as the connector edge.
The exact parameter estimates still give $|C_*|-K\to\infty$.
There are therefore enough other edges for the remaining binary detours.

Join one endpoint of $f$ to one endpoint of the designated edge $e_{\rm short}$
through all $B$ short blocks, in a fixed order.
Use $B+1$ additional $G_2$ paths of length $d+2$ between the fixed spanning-path endpoints.
Their interiors avoid every initial component and every earlier connecting path.
The total number of $G_2$ requests is $O(d)$.
Their total vertex cost is $O(d^2)$.

Here the corrected exposure argument remains valid for this larger number of requests.
For a first endpoint query, inspect $h_0=\lceil\beta n/100\rceil$ unused candidate edges.
Its failure probability is at most $(\ln n)^{-1}$.
The leaf query has at least $\beta n/80$ unused candidates and failure probability at most $(\ln n)^{-5/4}$.
Only the selected endpoint of $f$ appears in two requests.
In its first request, use it for the bounded first-endpoint query.
In its second request, use it for the leaf query.
All other endpoints appear once.
The total failure probability is
\[
 O\!\left(d\bigl((\ln n)^{-1}+(\ln n)^{-5/4}\bigr)\right)=o(1).
\]
The expansion and rooted-tree inputs are those proved above.

\emph{The additional paths in $G_3$.}
Embed the original binary detours and the two-digit detours with the audited procedure.
For every theta block, embed its additional path of length $z$ from the table.
This requires a rooted depth-$d$ tree with a stem of length $z-d-2\ge2$.
The new path uses $O(d)$ vertices, and all such paths use $O(d^2)$ vertices together.
Their endpoints lie in distinct base cycles and were fixed before $G_3$ was exposed.
The first-endpoint and leaf failure bounds above apply.
There are $O(d)$ such requests, so their total failure probability also tends to zero.
Expose only reservoir edges before querying these endpoint edges.

The total number of used vertices remains at most
\[
 L_0+O(d^2+b^3)=L_0+o(L_0)<2L_0.
\]
Thus the corrected robust expansion statement applies at every stage.

\emph{Hamilton completion and the edge count.}
Replace each short block by its fixed spanning path.
Their connecting chain joins one endpoint of $f$ to one endpoint of $e_{\rm short}$.
In the main binary unit, delete $f$ and replace every binary shortcut edge by its detour.
In the two-digit unit, delete $e_{\rm short}$ and replace every designated edge by its detour.
These replacements and the short-block chain form one spanning path of the exposed construction.
Its two endpoints are the unused endpoints of $f$ and $e_{\rm short}$.
Apply the unchanged $G_4$ Hamilton-path completion to these endpoints and the remaining vertices.

The resulting graph contains a Hamilton cycle.
Its edges outside that Hamilton cycle number
\[
 (K+1)+2+2b+(A+B)
 =K+2b+(T-2-4J)+3.
\]
Here a single-cycle block contributes one edge, and a theta block contributes two.
All short lengths $3,\ldots,T$ are present by Lemma~\ref{lem:rhost-theta-interval}.
The two-digit unit supplies $[T+1,T+b^2]$, which reaches $|C_*|$.
The binary unit supplies $[|C_*|,|C_*|+L_0]$ and $[n-L_0,n]$.
Its vertices remain on one contiguous interval of the Hamilton cycle.
The original long interval-covering argument in $G_5$ therefore applies without change.
It uses at most $m=\lfloor n2^{-K}\rfloor=O(\sqrt{\ln n})$ further chords.

The final excess is at most
\[
 K+2b+(T-2-4J)+m+3
 \le K+\frac32d+3b+m+8
 \le\log_2n+\frac32d+O(\sqrt{\ln n}).
\]
Since $d=(1+o(1))\ln n/\ln\ln n$, this proves the claim at $p_*$.
Monotonicity gives the result for $p\ge p_*$.
\end{proof}

Together, Theorems~\ref{thm:rhost-coefficient-one} and~\ref{thm:rhost-three-halves}
give, with high probability at $p=p_*$,
\[
 \boxed{\quad
 \log_2n+(1-o(1))\frac{\ln n}{\ln\ln n}
 \le\operatorname{Pex}(G(n,p_*))
 \le\log_2n+\left(\frac32+o(1)\right)\frac{\ln n}{\ln\ln n}.
 \quad}
\]

\subsection{Embedding whole theta blocks}
\label{subsec:rhost-whole-theta}

\begin{lemma}[Overlap estimate for a subdivided theta graph]
\label{lem:rhost-theta-moments}
Let $Q$ be a simple theta graph with $v$ vertices and $e=v+1$ edges.
Write $\ell_1,\ell_2,\ell_3$ for its three cycle lengths.
Let $J\sim G(n,D/n)$, with $D\ge54$ and $v^2/n=o(1)$.
Let $X$ count the copies of $Q$ in $J$.
Let $Y$ count ordered pairs of copies with a common vertex, including identical pairs.
Put
\[
 \mu=\mathbb E X=\frac{(n)_v}{|\operatorname{Aut}(Q)|}(D/n)^e,
 \qquad
 \varepsilon_Q=\frac{v^2}{n}
 +\sum_{i=1}^{3}\ell_iD^{-\ell_i}+nD^{-e}.
\]
There is an absolute constant $C$ such that
\[
 \operatorname{Var}(X)\le C\mu^2\varepsilon_Q,
 \qquad \mathbb E Y\le C\mu^2\varepsilon_Q.
\]
\end{lemma}

\begin{proof}
Use labelled embeddings first.
For two embeddings, their common vertices and common edges give a partially identified subgraph of $Q$.
An overlap with $u$ vertices and $f$ edges has normalized weight
\[
 (1+O(v^2/n))n^{-u}(n/D)^f.
\]
This follows by choosing the unshared vertices of the second embedding outside the first.
We sum over all possible overlap subgraphs and their injective maps into $Q$.
Allowing extra maps only increases the bound.

A connected tree overlap with $k$ edges has normalized weight $n^{-1}D^{-k}$.
There are at most $v9^k$ rooted tree subgraphs in the first copy.
A canonical depth-first traversal gives this bound, since the maximum degree is three.
There are at most $v3^k$ injective maps of each tree into the second copy.
The total contribution of one tree component is therefore at most
\[
 \frac{v^2}{n}\sum_{k\ge0}(27/D)^k\le\frac{2v^2}{n}.
\]
Allowing an unordered collection of tree components gives an exponential factor
$\exp(2v^2/n)$.
The nonempty forest overlaps contribute $O(v^2/n)$.

A proper connected overlap containing a cycle is unicyclic.
Its cycle is one of the three cycles of $Q$.
Outside that cycle, it contains only initial segments of the remaining theta path, from either endpoint.
For a cycle of length $\ell$, its image has at most three choices and at most $2\ell$ orientations and shifts.
Once this image is fixed, each additional initial segment has at most one possible image.
Summing their lengths gives a factor at most $(1-D^{-1})^{-2}$.
Thus these connected overlaps contribute at most
\[
 \frac{6}{(1-D^{-1})^2}\sum_{i=1}^{3}\ell_iD^{-\ell_i}.
\]
At most one unicyclic component can occur, since every two cycles of $Q$ meet.
Other components are trees and supply the same exponential factor as before.

An overlap with two independent cycles contains every edge of $Q$.
It contributes at most $12nD^{-e}$, because $|\operatorname{Aut}(Q)|\le12$.
These cases prove the bound for intersecting pairs of labelled embeddings.
Passing to unlabelled copies divides both pair counts and squared means by the same factor.

The contribution of vertex-disjoint pairs to $\mathbb E X^2$ is at most $\mu^2$.
Indeed, their vertex-placement factor is $(n-v)_v/(n)_v\le1$.
The remaining contribution is $\mathbb E Y$.
This also proves the variance bound.
\end{proof}

\begin{lemma}[Simultaneous disjoint theta embeddings]
\label{lem:rhost-theta-list}
Let $d\to\infty$, with $d=O(\ln n)$, and let $D=\Theta(d)$.
Assume $nD^{-d}$ tends to zero faster than every fixed negative power of $d$.
Consider $A=\Theta(d)$ simple theta graphs $Q_1,\ldots,Q_A$.
Suppose their numbers of edges lie between $d$ and $Cd$, for a fixed constant $C$.
Suppose their shortest cycle lengths are exactly $3,\ldots,A+2$,
and every other cycle has length at least $A+3$.
Then, with high probability, $G(n,D/n)$ contains vertex-disjoint copies of all $Q_i$.
They can also avoid any prescribed set of $O(d^2)$ vertices, even if that set is chosen after the graph is observed.
\end{lemma}

\begin{proof}
Choose $M=C_0d^2$, with $C_0$ larger than the total required vertex budget and the prescribed forbidden-set bound.
We show that every $Q_i$ has more than $M$ pairwise vertex-disjoint copies.
This permits a greedy choice of one copy of each type.

For one type, use $X,Y,\mu,\varepsilon_Q$ from Lemma~\ref{lem:rhost-theta-moments}.
Chebyshev's inequality bounds $\mathbb P(X<\mu/2)$ by $4C\varepsilon_Q$.
Markov's inequality bounds $\mathbb P(Y>\mu^2/(8M))$ by $8CM\varepsilon_Q$.
On the complementary event, form the graph whose vertices are the copies of $Q$.
Join two copies when they share a vertex.
It has an independent set of size at least $X^2/Y\ge2M$.
This bound follows by taking a random vertex order and retaining vertices earlier than all their neighbours.
Cauchy--Schwarz gives an expected retained size at least $X^2/Y$.

The sum of the error parameters over all types is
\[
 O\!\left(\frac{d^3}{n}+D^{-3}+d\,nD^{-d}\right).
\]
Here the distinct short cycle lengths give
$\sum_{s\ge3}sD^{-s}=O(D^{-3})$.
The longer cycle lengths give a smaller contribution.
After multiplication by $M=O(d^2)$, this sum tends to zero.
Indeed, $d^5/n=o(1)$ and $d^2D^{-3}=O(1/d)$.
The hypothesis handles the final term.

Each forbidden vertex can meet at most one member of a vertex-disjoint family of copies.
At every greedy step, fewer than $M$ vertices have been excluded.
An available copy therefore remains.
\end{proof}

\begin{lemma}[A denser packing of the required lengths]
\label{lem:rhost-whole-theta-interval}
Put $T=2(d+b+1)$, where $b\to\infty$ and $b=o(d)$.
Set
\[
 J=\left\lfloor\frac{T-2}{12}\right\rfloor,\qquad A=4J.
\]
There are $A$ theta blocks covering exactly the lengths $3,\ldots,3A+2$.
Their shortest cycle lengths are $3,\ldots,A+2$.
Their other cycle lengths are at least $A+3$.
Each block has between $d$ and $O(d)$ edges.
With one single-cycle block for each remaining length through $T$,
the total rank is at most
\[
 \frac43d+\frac43b+4.
\]
\end{lemma}

\begin{proof}
For $0\le j<J$, put $q=J-j-1$ and $h=A+3+8j$.
Use the four triples
\[
 (4q+3,h+4,h+5),\quad (4q+4,h,h+2),\quad
 (4q+5,h+6,h+7),\quad (4q+6,h+1,h+3).
\]
Lemma~\ref{lem:rhost-theta-block} gives the required theta graphs.
The low triples run through $3,\ldots,A+2$ in reverse blocks.
The high triples run through every length $A+3,\ldots,3A+2$ once.
The four edge counts in block $j$ are
\[
 6J+6j+7,\quad6J+6j+4,\quad6J+6j+10,\quad6J+6j+6.
\]
Their minimum is at least $6J+4\ge d+b-1\ge d$.
Their maximum is $O(d)$.
There are at most eleven remaining lengths through $T$.
The total rank is
\[
 2A+(T-2-3A)=T-2-A
 \le\frac23(T-2)+4=\frac43d+\frac43b+4.
\]
\end{proof}

\begin{theorem}[An upper coefficient of four thirds]
\label{thm:rhost-four-thirds}
For the same $p_*$ as in Theorem~\ref{thm:rhost-scale}, and every $p\ge p_*$,
with high probability
\[
 \operatorname{Pex}(G(n,p))
 \le\log_2n+\left(\frac43+o(1)\right)\frac{\ln n}{\ln\ln n}.
\]
\end{theorem}

\begin{proof}
Split $G_2$ into two independent subexposures of probability
$q_2=1-\sqrt{1-p_2}$.
Their union has the original probability $p_2$.
Replace the expansion parameter by
\[
 \beta=\frac{4(\ln\ln n)^2}{\ln n},\qquad
 a=\left\lfloor(5\beta)^{-1}\right\rfloor,
 \qquad d=\left\lfloor\log_a(n/200)\right\rfloor.
\]
The remaining parameters retain their previous values.
The depth satisfies $d=(1+o(1))\ln n/\ln\ln n$, and the second $G_2$ subexposure obeys
\[
 q_2\beta n\ge100\ln\ln n.
\]
Thus the corrected expansion and endpoint-query bounds retain their stated constants.
They also hold in $G_3$, whose edge probability is larger.

Prepare the binary base cycles, the two-digit base cycle, the reserve edge $f$,
and the at most eleven single-cycle blocks in $G_1$, as above.
Use the first $G_2$ subexposure to embed the complete theta list from
Lemma~\ref{lem:rhost-whole-theta-interval}, avoiding these initial vertices.
This is justified by Lemma~\ref{lem:rhost-theta-list}.
Indeed, $D=nq_2=\Theta(d)$, and
\[
 \ln D-\ln a\ge\ln(500\ln\ln n).
\]
Consequently,
\[
 nD^{-d}
 \le\exp\!\left(-\Omega\!\left(
 \frac{\ln n\,\ln\ln\ln n}{\ln\ln n}\right)\right).
\]
This is smaller than every fixed negative power of $d$.

All short blocks now exist before the second $G_2$ subexposure.
Choose their spanning-path endpoints and join them into the same connector chain as before.
Construct the binary base cycle through $f$ in this second subexposure as well.
There are $O(d)$ endpoint requests, with total failure probability $o(1)$.
Embed the binary and two-digit detours in $G_3$.
No additional theta path is needed there.
The vertex budget is still $L_0+O(d^2+b^3)<2L_0$.

Use the same $G_4$ Hamilton completion and $G_5$ long interval-covering chords.
The short blocks cover $3,\ldots,T$.
Their total rank is at most $4d/3+4b/3+4$.
Thus the excess is at most
\[
 K+2b+\left(\frac43d+\frac43b+4\right)+m+3
 \le\log_2n+\frac43d+O(\sqrt{\ln n}).
\]
The five original exposure probabilities have not increased.
This proves the claim at $p_*$, and monotonicity gives the remaining range.
\end{proof}

Combining the coefficient-one lower bound with the whole-theta construction gives
\[
 \boxed{\quad
 \log_2n+(1-o(1))\frac{\ln n}{\ln\ln n}
 \le\operatorname{Pex}(G(n,p_*))
 \le\log_2n+\left(\frac43+o(1)\right)\frac{\ln n}{\ln\ln n}
 \quad}
\]
with high probability.
This construction leaves the sharp coefficient and the full-threshold question unresolved.
Both are treated further in the later sections.
